座標の変遷は次のようになる。
\[
\xymatrix@R=40pt{
(1,-3,1)\ar[d]^{表}&\\
(-1,2,-1)\ar[d]^{表}&\\
\ar@/_10pt/[d]_{裏} (1,-1,0)\ar@/_10pt/[r]_{表}&\ar@/_10pt/[l]_{表}(-1,1,0)\ar@/_10pt/[d]_{裏}\\
\ar@/_10pt/[u]_{裏} (0,-1,1)\ar@/_10pt/[rd]_{表}& \ar@/_10pt/[u]_{裏} (0,1,-1)\ar[d]^{表}\\
&(0,0,0)
}
\]
そこで, $n\geq 0$に対して次のように数式を定める:
\begin{align*}
a_n &= ( \mr{P}_{n} = (1,-3,1) である確率)\\
b_n &= ( \mr{P}_{n} = (-1,2,-1) である確率)\\
c_n &= ( \mr{P}_{n} = (1,-1,0)または(-1,1,0) である確率)\\
d_n &= ( \mr{P}_{n} = (0,-1,1)または(0,1,-1) である確率)\\
e_n &= (\mr{P}_{n} = (0,0,0)である確率)
\end{align*}
初期状態より$a_0=1$, $b_0=c_0=d_0=e_0 = 0$であり, 全事象の確率から$a_n + b_n + c_n + d_n + e_n = 1$である。
変遷図から, $n\geq 0$において次の漸化式が得られる。
\[
\bcas
a_{n+1} = \dfrac{1}{2}a_n \\
b_{n+1} = \dfrac{a_n+b_n}{2} \\
c_{n+1} = \dfrac{b_n + c_n + d_n}{2}\\
d_{n+1} = \dfrac{c_n}{2}\\
e_{n+1} = \dfrac{d_n}{2} + e_n
\ecas
\]
求めるべきは$e_n$ ($n\geq 1$)の一般項である。まず, 上から簡単に分かるのが
\[a_{n} = \dfrac{1}{2^{n}}\]
\[c_{n+2} = \dfrac{b_{n+1}}{2} + \dfrac{c_{n+1}}{2} + \dfrac{c_n}{4}  \]
\[\bolm{e_n = 1-a_n - (b_n + c_n + d_n) = 1-\dfrac{1}{2^{n}} - 2c_{n+1}}\]
である。したがって, $b_n$を決定すれば二式目より$c_n$の漸化式を解くことができ, $c_n$を決定すれば三式目より$e_n$が決定される。\par
$(0,0,0)$になるにはコイントスを少なくとも4回はしなければならないので$e_0=e_1=e_2=e_3=0$も分かる。\par
$a_n$の結果を用いて$b_{n}$を求めると,
\begin{align*}
b_{n+1} &= \dfrac{b_n}{2} + \dfrac{1}{2^{n+1}}\\
&\iff 2^{n+1}b_{n+1} - 2^{n}b_n = 1
\end{align*}
であるから, 数列$\{ 2^{n} b_n\}$は第0項が$0$の階差1の数列なので$2^{n}b_n = n$である。よって$b_n = \dfrac{n}{2^{n}}$である。\par
次に$c_{n}$を計算する。$b_n$の結果から
\[c_{n+2} = \dfrac{n+1}{2^{n+2}} + \dfrac{c_{n+1}}{2} + \dfrac{c_n}{4}\quad (n\geq 0)\]
である。$2^{n}c_n = f_n$とおくと
\[f_{n+2} = (n+1) + f_{n+1} + f_n\quad (n\geq 0)\]
である。変形すると
\[f_{n+2} + n+4 = (f_{n+1} + n+3) + (f_n + n+2)\]
であるから\footnote{$f_n + An + C$に関する非斉次な3項間漸化式になると予想した上で定数$A,C$の満たすべき条件を見れば良い。}, $g_{n} = f_n + n+2$とおくと
\[g_{n+2} = g_{n+1} + g_n\quad (n\geq 0)\]
となる。特性方程式の解を見ることで$g_{n} = A_1\parena{\dfrac{1+\sqrt{5}}{2}}^{n} + A_2\parena{\dfrac{1-\sqrt{5}}{2}}^{n}$の形と分かる。$g_0 = 2, g_1 = 3$であるから$A_1 = 1+\dfrac{2}{\sqrt{5}}$, $A_2 = 1-\dfrac{2}{\sqrt{5}}$となる。簡単のため$\phi = \dfrac{1+\sqrt{5}}{2}$, $\psi=\dfrac{1-\sqrt{5}}{2}$とおくと, $n\geq 0$において
\begin{align*}
c_n &= \dfrac{g_n - n -2}{2^{n}} \\
&= A_1\parena{\dfrac{\phi}{2}}^{n} + A_2\parena{\dfrac{\psi}{2}}^{n}  - \dfrac{n+2}{2^{n}}
\end{align*}
よって, $n\geq 0$において
{\footnotesize
\begin{align*}
e_n &= 1 - \dfrac{1}{2^{n}} - 2A_1\parena{\dfrac{\phi}{2}}^{n+1} - 2A_2\parena{\dfrac{\psi}{2}}^{n+1} + \dfrac{n+3}{2^{n}}\\
&=\textcolor{blue}{ 1+ \dfrac{n+2}{2^{n}} - \dfrac{10+4\sqrt{5}}{5}\parena{\dfrac{1+\sqrt{5}}{4}}^{n+1} - \dfrac{10-4\sqrt{5}}{5}\parena{\dfrac{1-\sqrt{5}}{4}}^{n+1}}
\end{align*}
}
